Training and serving a modern Large Language Model stresses every
layer of the deep-learning software stack at once: the distributed
schedule must keep thousands of accelerators coordinated, the kernel
layer must drive Tensor Cores near peak, the attention algorithm must
scale sub-quadratically with the context window, and the serving
runtime must turn the accumulated algorithmic savings into end-user
throughput. Isolated optimisations at any one of these layers leave
the gains of the others on the table. This thesis develops, implements,
and evaluates an end-to-end stack in which the five layers are
designed to compose, and measures its behaviour on NVIDIA A100 hardware
across both training and inference workloads.

The work addresses LLM optimisation along five complementary axes.
\emph{First}, a systematic comparative study of four distributed training
strategies---Data Parallelism (DP), Distributed Data Parallel (DDP), Fully
Sharded Data Parallel (FSDP), and DeepSpeed ZeRO-3---is carried out on a
single-node 8-GPU A100 server with identical workloads, identical
initialisations, and identical random seeds, isolating the effect of the
parallelisation strategy itself. \emph{Second}, low-level matrix-multiplication
kernels are built using NVIDIA CUTLASS 3.x and the CuTe layout algebra,
covering threadblock/warp/instruction tiling, multi-stage software pipelining,
Tensor-Core utilisation, and epilogue fusion for bias, activation, and
residual additions. \emph{Third}, a novel block-sparse \emph{Sliding Window
Flash Attention} (SWFA) algorithm is developed that augments Flash Attention
with an explicit block-classification step and a vectorised mask generator.
Entire $B_r\!\times\!B_c$ blocks that fall outside the attention band are
skipped rather than merely masked. \emph{Fourth}, GPU kernels inside the
SGLang serving runtime are engineered for long-context inference, including
paged KV cache access, fused decode attention, and fused top-$k$/top-$p$
sampling. \emph{Fifth}, a Triton-based kernel profiling and auto-tuning
pipeline is assembled that combines CUDA-event timing, roofline analysis,
Nsight Compute cross-checks, and shape-indexed configuration caching.

Experimental results on an 8-GPU A100-40GB server show that DeepSpeed ZeRO-3
attains the highest GPU utilisation ($93.8\%$ peak, $65.8\%$ average) and the
best memory efficiency ($\approx 87.5\%$ reduction per GPU over DP), while
FSDP provides the best memory/throughput trade-off for medium-scale models.
The proposed SWFA delivers up to $65.6\%$ reduction in computed attention
blocks for narrow windows while matching the reference output within
$10^{-4}$ FP32 tolerance and $10^{-2}$ FP16 tolerance across nineteen
validation configurations. When combined with DeepSpeed or FSDP on a
GPT-2-Medium workload with sequence length $2048$, SWFA enables end-to-end
training throughput improvements of $50$--$59\%$ and energy savings of
$25$--$29\%$ per epoch. A closed-form expression for the optimal window
size under an exponential attention-decay cost model is derived, and strict
convexity of the corresponding loss is proved. The SGLang kernel layer
achieves a cumulative $2.84\times$ decode-throughput speedup on LLaMA-7B,
and the Triton profiling pipeline reproduces these numbers independently
with $<\!3\%$ deviation.

\vspace{4mm}
\noindent\textbf{Keywords:} Large Language Models, Distributed Training,
Flash Attention, Sliding-Window Attention, CUTLASS, CuTe, SGLang,
Triton, GPU Computing, Kernel Profiling, Tensor Cores.
